\chapter{Evaluation Protocol}
\label{cha:protocol}

This chapter is short and load-bearing. Every number in Chapters~\ref{cha:baselines}
and~\ref{cha:expdesign} is only as credible as the protocol that produced it, and the
protocol is where a multimodal cross-plant study is most easily compromised --- by a
normaliser fitted on test plants, by a horizon that differs between models, by a retrieval
datastore that quietly contains the answer.

\section{Three fairness principles}
\label{sec:fairness}

\begin{enumerate}
  \item \textbf{Same horizon and granularity for every model.} All models forecast the same
    future window from the same history cadence. Context length is itself a fairness
    variable: a model given more history is not a better model, so no comparison is made
    across differing $T$.
  \item \textbf{Disjoint test plants, with no statistic leakage.} No model may train on, or
    derive statistics from, held-out plants. This includes normalisers --- hence
    capacity normalisation (Section~\ref{sec:harmonisation}) rather than per-series
    standardisation.
  \item \textbf{No domain physics inside models.} Clear-sky-index conversions and
    irradiance-to-power formulas are excluded from every learned model, so that measured
    improvements are attributable to modelling rather than to a hand-coded solar prior. The
    sole exemption is smart persistence, which is the physics reference itself.
\end{enumerate}

\subsection{The failure modes these principles prevent}

Each principle exists because of a specific way a cross-plant multimodal benchmark can be
compromised, usually without anyone noticing. Naming the failure mode makes the rule
checkable.

\paragraph{Normaliser leakage.} The most common and least visible failure. Standardising each
series by its own mean and standard deviation is routine in forecasting, and under a
cross-plant protocol it silently passes test-plant statistics into the model. The symptom is a
benchmark on which every method looks unusually strong and no method generalises in
deployment. Capacity normalisation eliminates the channel entirely, because nameplate capacity
is metadata, not a measurement.

\paragraph{Context-length mismatch.} If one model receives fourteen days and another six
hours, the comparison measures history length rather than architecture. The failure is easy to
introduce accidentally, because foundation models have fixed maximum contexts and quietly
truncate. The rule here is that context length is a fairness variable: models that cannot
consume the full window are documented as such rather than silently compared.

\paragraph{Transductive retrieval.} A retrieval baseline whose datastore contains test-plant
history is performing in-context memorisation of the evaluation target. The resulting number
can be excellent and means nothing about zero-shot transfer. The datastore rule closes this,
and the transductive condition --- which is a legitimate thing to study --- is reported
separately if at all.

\paragraph{Privileged covariates.} Supplying observed future weather over the horizon is not
forecasting; it is forecasting given a perfect weather forecast. Because covariate arrays
naturally span history and horizon, this is a single indexing decision away at all times. The
protocol partitions covariates into deterministic and observed classes
(Appendix~\ref{app:covariates}), and models receiving the observed class over the horizon are
labelled oracle and excluded from competitive comparison.

\paragraph{Physics smuggled into a model.} A clear-sky conversion inside a learned model
converts a modelling result into a statement about how good the clear-sky model is. Since
smart persistence --- which is that physics --- is the reference, a model that internalises it
is partly competing against itself. The exemption is granted to exactly one model, and it is
the reference.

\paragraph{Metric-window mismatch.} If two models are scored on different evaluation windows,
their metrics are not comparable even under an identical protocol. This is the defect behind
the flagged Time-VLM row of Chapter~\ref{cha:baselines}, and it is the reason skill scores are
always computed against the reference evaluated on the same windows as the model being
scored, rather than against a single global reference value.

\section{Windows in physical time}
\label{sec:windows}

Windows are defined in physical time and resolved to steps per dataset, rather than being
defined in steps.

\begin{itemize}
  \item \textbf{History $T$ = 14 days} --- 672 steps at the \texttt{uk\_pv} 30-minute
    cadence, 1344 at the \texttt{goes\_pvdaq} 15-minute cadence. Two weeks is
    deployment-realistic (a newly connected plant will have that much history) and gives
    every model multiple diurnal cycles, which shorter contexts denied the foundation
    models in preliminary work. Foundation models with a fixed maximum context consume
    their most recent supported window.
  \item \textbf{Horizon $H$ = 6 hours} --- 12 and 24 steps respectively. Skill-decay curves
    are additionally reported at 1, 6 and 24 hours, the last covering the day-ahead case.
  \item \textbf{Visual window $T_v$ = 8 frames} over a short recent window of three to six
    hours, \emph{decoupled} from the numerical history. Vision carries information only over
    the cloud-advection horizon; widening the visual window past a few hours adds compute,
    not signal. This asymmetry --- long numerical context, short visual context --- is a
    design commitment of the architecture in Chapter~\ref{cha:method}, not merely a budget
    decision.
\end{itemize}

\paragraph{The cadence rule.} Because windows are physical, step counts differ across
datasets: 672 against 1344 for the same fourteen days. A step-horizon of 12 therefore means
six hours in one dataset and three in the other. Two hard rules follow. Physical lead time
is reported beside every per-dataset table; and metrics are never pooled across cadences in
raw form --- cross-dataset aggregation uses only the scale-free statistics of
Section~\ref{sec:aggregation}.

\section{Metrics}
\label{sec:metrics}

All metrics are computed on the capacity-normalised target over daylight, valid steps, per
plant, and then macro-averaged over plants so that large or data-rich plants do not
dominate.

\paragraph{Point accuracy.} With $M$ evaluation samples, horizon $H$ and plant capacity
$C_i$:
\[
  \mathrm{NMAE} = \frac{1}{M H} \sum_{i=1}^{M} \sum_{h=1}^{H}
  \frac{|\hat{y}_{i,h} - y_{i,h}|}{C_i},
  \qquad
  \mathrm{NRMSE} = \sqrt{\frac{1}{M H} \sum_{i=1}^{M} \sum_{h=1}^{H}
  \left( \frac{\hat{y}_{i,h} - y_{i,h}}{C_i} \right)^{\!2}}.
\]

\paragraph{Skill score.} The headline number, defined against smart persistence and
computed from NRMSE:
\[
  \mathrm{SS} = 1 - \frac{\mathrm{NRMSE}_{\text{model}}}
                        {\mathrm{NRMSE}_{\text{smart persistence}}}.
\]
A score of zero means no better than the reference; one means perfect; negative means
worse than a model that assumes cloud conditions persist. An NMAE-based variant may be
reported as a secondary column but must be labelled as such, since the two rank models
differently.

\paragraph{Probabilistic accuracy.} For models with a distributional output, the continuous
ranked probability score is approximated by the pinball loss averaged over the quantile
grid~\cite{gneiting}. It rewards forecasts that are both accurate and honestly calibrated,
and it is the only metric here that penalises overconfidence. Calibration is additionally
reported as empirical coverage of the nominal 80\% interval.

\paragraph{Variance tracking.} The squared Pearson correlation $R^2$ between prediction and
truth measures how much of the \emph{shape} of the day a forecast tracks, independently of
bias or scale.

\paragraph{Reading $R^2$ and skill together.} These two must be read jointly, because their
disagreement is diagnostic. A model with high $R^2$ and low skill tracks the timing of the
day correctly but sits off the diagonal --- a bias or scale problem, fixable by
recalibration. A model with low $R^2$ and moderate skill is doing the opposite. In
Chapter~\ref{cha:baselines} several multimodal baselines exhibit exactly the first pattern,
and the distinction changes the interpretation of their ranking entirely.

\paragraph{Ramp regime.} Point metrics are additionally reported on a \emph{ramp subset}:
the top decile of $|\Delta y|$ per site, that is, the cloud-transition periods. This is
where a visual modality should earn its tokens --- clear-sky periods are won by persistence
--- so the ramp columns, not the aggregate ones, are the sharpest test of the multimodal
hypothesis.

\subsection{Metrics considered and rejected}

Three alternatives were considered and are recorded here with the reason for rejection,
because a metric choice that is not justified is a metric choice a reader has to take on
trust.

\paragraph{Mean absolute percentage error.} Undefined at zero and unbounded near it. The
target is exactly zero for roughly half of every day and arbitrarily close to it at dawn and
dusk, so any percentage-of-actual metric is dominated by the least interesting observations
in the dataset.

\paragraph{Per-series standardised error.} Normalising each plant's series by its own mean and
standard deviation is standard in general-purpose forecasting benchmarks. It is prohibited
here: those statistics would have to be computed on test-plant data, which is precisely the
leakage the cross-plant protocol exists to prevent. Capacity normalisation uses a quantity
known at installation and is leak-free by construction.

\paragraph{Energy-weighted error.} Weighting errors by absolute energy would emphasise
high-output periods, which is arguably what a grid operator cares about. It was rejected
because it would make the metric depend on the capacity mix of the test set, undermining
comparability across the two datasets --- and because the ramp subset already isolates the
high-consequence regime in a way that is transparent rather than baked into the headline
number.

\subsection{Aggregation order matters}

Metrics are computed \emph{per plant and then averaged}, not pooled over all observations.
The two differ whenever plants contribute unequal numbers of valid steps or have different
capacity factors, which they do (Figure~\ref{fig:siteboxplot}, Figure~\ref{fig:availability}).
Pooling would weight a data-rich, high-output plant more heavily than a sparse one, and the
resulting number would answer ``how well does the model do on the average \emph{observation}''
rather than ``how well does it transfer to the average \emph{plant}''. Since the research
question is about generalisation across plants, each plant is one unit of evidence.

The choice has a practical consequence worth noting: macro-averaged metrics are noisier than
pooled ones, because a single difficult plant is not diluted. With fourteen test plants the
standard error of the mean is not negligible, which is a further reason to treat small
differences between models as ties.

\section{Scale-free aggregation}
\label{sec:aggregation}

When results must be summarised across datasets of different cadence, or across scenarios,
only rank-preserving and scale-free statistics are used: win rate against the reference,
geometric-mean skill, and average rank. Raw step-horizon metrics are never pooled across
cadences.

\section{Scenario battery}
\label{sec:scenarios}

Seven scenarios were specified. Table~\ref{tab:scenarios} states each and, in the last
column, whether it was executed for this thesis --- an honest accounting matters more than
a complete-looking table.

\begin{table}[h!]
\centering
\small
\begin{tabular}{l p{5.4cm} p{4.4cm} l}
\hline
\textbf{ID} & \textbf{Split} & \textbf{Question} & \textbf{Run} \\
\hline
S1 & Train plants, held-out time range & Sanity / upper bound & no \\
S2 & \textbf{Disjoint test plants (primary)} & Headline result (H3) & \textbf{yes} \\
S3 & Train \texttt{uk\_pv}, test \texttt{goes\_pvdaq} and reverse & Distribution-shift
generalisation & no \\
S4 & S2 with $H \in \{1,6,24\}$\,h & Skill-decay curves & partial \\
S5 & S2 with 10/25/50/100\% of train plants & Sample efficiency & no \\
S6 & S2 restricted to top-decile $|\Delta y|$ windows & Where vision should win &
\textbf{yes} \\
S7 & Train on a month subset, test unseen season & Temporal shift & no \\
\hline
\end{tabular}
\caption{Scenario battery. Everything reported in this thesis is S2 on \texttt{uk\_pv},
plus the S6 ramp subset and a partial S4 decay profile.}
\label{tab:scenarios}
\end{table}

\subsection{What each scenario is for}

\paragraph{S1, in-domain.} Train plants, held-out time range. This is the protocol most of the
literature reports, included here only as a sanity ceiling: the gap between S1 and S2 measures
how much of a model's apparent accuracy is plant-specific memorisation.

\paragraph{S2, cross-plant.} The headline. Disjoint test plants, short inference history, no
test-plant statistic anywhere. Everything reported in this thesis is S2.

\paragraph{S3, cross-dataset.} Train on one track, test on the other. This is the strongest
generalisation test available in the data, because it shifts sensor, scale, climate, latitude
and cadence at once (Section~\ref{sec:eda}). It also involves a cadence change, so results
must be reported with the physical lead time attached and aggregated only through the
scale-free statistics of Section~\ref{sec:aggregation}. A robustness variant resamples the
faster track to the slower cadence so that step-horizon and physical horizon coincide.

\paragraph{S4, long horizon.} The same protocol at one, six and twenty-four hours. Skill decay
curves are how a forecaster's useful range is communicated, and they separate models that are
good at short lead times from models that are good at persistence.

\paragraph{S5, data efficiency.} S2 with 10, 25, 50 and 100 percent of training plants. This
is the canonical foundation-model claim --- that pretraining buys sample efficiency --- and it
is the scenario in which a pretrained backbone should beat a supervised one even if it does
not at full data. Its absence is why the comparison in Chapter~\ref{cha:benchmark} between
supervised and pretrained models must be read as holding \emph{at 69 training plants} and not
in general.

\paragraph{S6, ramp subset.} S2 restricted to the top decile of $|\Delta y|$. The regime where
a visual modality should win, and therefore the sharpest test of the thesis.

\paragraph{S7, seasonal transfer.} Train on a subset of months, test on an unseen season. A
temporal analogue of the cross-plant shift, included for completeness and not executed.

The unexecuted scenarios are not decoration. S3 and S5 are the standard evidence a
foundation-model claim is expected to provide --- transfer under distribution shift, and
accuracy as a function of training data --- and their absence is the principal limitation
of this work, restated in Section~\ref{sec:limitations}.

\section{Input parity}
\label{sec:parity}

Table~\ref{tab:parity} states which inputs each tier may consume. A baseline never receives
information the models it is compared against cannot access.

\begin{table}[h!]
\centering
\small
\begin{tabular}{l c c c c c}
\hline
\textbf{Tier} & $\mathbf{Y}$ & $\mathbf{X}_{\text{cov}}$ & $\mathbf{V}$ &
\textbf{Retrieval} & \textbf{Text} \\
\hline
T0 Reference & yes & clear-sky only & --- & --- & --- \\
T1/T2 Supervised & yes & yes & --- & --- & --- \\
T3 Foundation models & yes & where native & --- & --- & --- \\
T4 Adaptation & yes & yes (CoRA) & --- & yes (RAG) & --- \\
T5/T6 Multimodal & yes & yes & yes & model-specific & model-specific \\
MMTSFM (ours) & yes & yes & yes & optional (H4) & --- \\
\hline
\end{tabular}
\caption{Input-parity matrix.}
\label{tab:parity}
\end{table}

Two rules within this matrix are the ones a critical reader attacks first, and both are
enforced.

\paragraph{The retrieval datastore rule.} Retrieval and memory baselines may populate their
datastore with \emph{training-plant data only}. Allowing test-plant history into the
datastore converts zero-shot forecasting into transduction, and would make the retrieval
tier incomparable with everything else. Transductive retrieval is a legitimate but
\emph{separate} condition and never enters the headline table.

\paragraph{The oracle-covariate rule.} Two Chronos-2 variants are reported that consume
\emph{observed future weather} over the forecast horizon --- temperature, cloud cover,
irradiance that would not be known at forecast time. They are labelled \emph{oracle} and
are upper bounds, not competitors. The honest variants receive only deterministic future
covariates: solar geometry and calendar encodings. The gap between the two is one of the
most informative measurements in this thesis (Section~\ref{sec:bar}), but it must
never be read as a competitive result.

Two further parity constraints apply to the multimodal tier. Models that natively want text
(Solar-VLM) may generate weather descriptions only from covariates available to every other
model --- no external weather service. And no model may consume a normaliser, calibration
constant or context-selection heuristic derived from the test plants.

\section{Reproducibility}
\label{sec:reproducibility}

Reproducibility here means something stronger than ``the code is available''. It means that a
number in this thesis can be traced to the job that produced it, that the job can be
re-specified exactly, and that the data it consumed cannot have changed underneath it. Three
mechanisms enforce that.

\emph{The data is immutable.} The dataset of record is treated as read-only, and no result
depends on re-processing it. Splits are committed to version control rather than regenerated
at run time, so a later run cannot silently receive a different partition; their disjointness
is asserted at load, so a leak raises rather than quietly improving a score.

\emph{The configuration is declarative and complete.} Every hyperparameter is composed by
Hydra from configuration groups, and none is hardcoded in model code. Each baseline's
configuration is self-contained in its own directory, so no model's behaviour depends on a
global setting that another model also reads --- a failure mode that is invisible until two
models disagree for reasons unrelated to their architectures.

\emph{The experiment is registered before it runs.} Each experiment declares a one-sentence
hypothesis, a configuration diff, a branch and a nominated comparison baseline before the job
is submitted. This is a guard against a specific and common failure: reframing an experiment
after the fact as a test of whatever it happened to show.

A fixed global seed of 42 governs data loaders, weight initialisation and dropout.
Configuration is Hydra-only~\cite{hydra}, with each baseline's hyperparameters
self-contained in its own directory, so that no model's behaviour depends on a global
setting another model also reads. Dependencies are managed with a single lockfile. Training
runs on the Leonardo cluster under SLURM, with standard output and error retained per job,
so that any reported number can be traced to the job that produced it. Splits are committed
to version control rather than regenerated, and their disjointness is asserted at load time.

\par

